\subsection{Cosmological Consequences of a Rolling Homogeneous Scalar Field --- 1988}

\textbf{Authors:} Bharat Ratra, P. J. E. Peebles

\paragraph{主な主張}
時間発展する一様な scalar field を宇宙成分として導入し、減少する nonlinear potential のもとで scalar energy density が特徴的に redshift する解と attractive time-dependent fixed point を示した \cite{Ratra:1988rolling}。

\paragraph{新規性}
後の quintessence/tracker 研究につながる rolling scalar cosmology を早期に具体化し、inverse-power-law 型 potential を含む動的 dark-energy model の基礎を与えた。

\paragraph{位置付け}
Inverse-power-law tracker の先駆的原論文として位置付けられ、Zlatev--Wang--Steinhardt の tracker 概念より前の重要な precursor である。

\paragraph{コメント}
\textbf{何の原論文か:} \textbf{inverse-power-law 型 rolling scalar / tracker-like quintessence の先駆的原論文}。Tracker という名称自体の原論文は Zlatev--Wang--Steinhardt (1999) である。
